% Abstract, revised 18 Nov 2026 from J.M.'s suggested text. Results come from the draft Chapter 4. Must fit one page.
The red-striped soft scale insect (RSSI), \emph{Pulvinaria tenuivalvata}, had affected 72,103 ha of sugarcane on Negros Island by 2 October 2026, 30.1\% of the area planted \citep{SRA2026FOI}. Its effect on yield has not been measured in the Philippines \citep{Guerrero2024}, yet Negros Occidental funds biological control \citep{Bayoran2026} while Negros Oriental planned chemical spraying \citep{Partlow2026}. A systematic review of 30 agent-based pest models found none of a scale insect and none coupling a farmer's choice between biological and chemical control with pest biology.

This study develops an agent-based model of RSSI management in GAMA. Pest development follows Philippine laboratory biology \citep{Guerrero2026} on a 1-ha map of Negros cane \citep{DiTommaso2024}. Over three crop years, under one common treatment trigger, the model compares \emph{no control}, \emph{chemical control} whose efficacy may decline with repeated use, fungal \emph{biological control}, and a \emph{hybrid} that applies the fungus at infestation Levels 1 and 2 and a chemical at Level 3.

Every control strategy cut cumulative pest load at least 82\% below no control. Chemical control matched biological control only if its efficacy halved no faster than about once per 80 applications. The hybrid added nothing while fungal efficacy held, but was best when both agents lost efficacy quickly. The ranking held in 15 of 18 sensitivity cases. A hindcast against SRA data showed baseline spread is too slow. These are comparisons between strategies under stated assumptions, not field predictions.
